\DocumentMetadata{
  pdfversion=2.0,
  pdfstandard=ua-2,
  lang=en-US,
  tagging=on,
  tagging-setup={math/setup=mathml-SE,tabsorder=structure}
}
\documentclass[11pt,letterpaper]{article}
\usepackage[margin=0.68in,includefoot]{geometry}
\usepackage{newcomputermodern}
\usepackage{amsmath,amscd,mathtools}
\usepackage{tikz,adjustbox}
\providecommand{\square}{\mdlgwhtsquare}
\usepackage[hidelinks]{hyperref}
\hypersetup{
  pdftitle={Algebra First-Year Exam, May 2014, Part 2},
  pdfauthor={University of Florida Department of Mathematics},
  pdfsubject={Graduate mathematics examination},
  pdfdisplaydoctitle=true,
  bookmarksopen=true
}
\setcounter{secnumdepth}{0}
\setlength{\parindent}{0pt}
\setlength{\parskip}{0.28em}
\setlength{\emergencystretch}{3em}
\setlength{\leftmargini}{2.15em}
\setlength{\leftmarginii}{2.65em}
\setlength{\leftmarginiii}{3em}
\pagestyle{empty}
\raggedbottom
\allowdisplaybreaks
\NewTaggingSocketPlug{block/list/label}{examproblem}{%
  \tagstructbegin{tag=Lbl}\tagstructend%
  \tagstructbegin{tag=\LBody}%
  \tagstructbegin{tag=\ProblemHeadingTag,title-o={\ProblemHeadingText}}%
  \tagmcbegin{tag=\ProblemHeadingTag,actualtext={\ProblemHeadingText}}%
  #2%
  \tagmcend%
  \tagstructend%
}
\newenvironment{examproblems}[2]{%
  \def\ProblemHeadingTag{#2}%
  \AssignTaggingSocketPlug{block/list/label}{examproblem}%
  \begin{enumerate}%
  \setcounter{enumi}{#1}%
  \renewcommand{\labelenumi}{\textbf{\theenumi.}}%
  \setlength{\topsep}{0.35em}%
  \setlength{\partopsep}{0pt}%
  \setlength{\itemsep}{0.48em}%
  \setlength{\parsep}{0.18em}%
}{\end{enumerate}}
\newenvironment{examsubparts}{%
  \AssignTaggingSocketPlug{block/list/label}{default}%
  \begin{enumerate}%
  \setlength{\topsep}{0.18em}%
  \setlength{\partopsep}{0pt}%
  \setlength{\itemsep}{0.16em}%
  \setlength{\parsep}{0.1em}%
}{\end{enumerate}}
\newenvironment{examinstructions}{%
  \par\begingroup\itshape
}{%
  \par\endgroup\vspace{0.18em}
}
\makeatletter
\renewcommand\subsection{\@startsection{subsection}{2}{0pt}%
  {-1.25ex plus -.25ex minus -.1ex}%
  {0.55ex plus .1ex}%
  {\normalfont\large\bfseries}}
\renewcommand{\@maketitle}{%
  \newpage\null\vskip -1em%
  {\footnotesize\bfseries
    \href{https://www.ufl.edu/}{University of Florida}, \href{https://math.ufl.edu/}{Department of Mathematics}\par}%
  \vskip -0.5em%
  \tagstructbegin{tag=H1,title={Algebra First-Year Exam, May 2014, Part 2}}%
  \tagpdfparaOff%
  \tagmcbegin{tag=H1}%
    {\LARGE\bfseries\href{https://gma.math.ufl.edu/exams/algebra-first-year/}{Algebra First-Year Exam}, May 2014, Part 2\par}%
  \tagmcend%
  \tagpdfparaOn%
  \tagstructend%
  \vskip 0.25em%
  \tagmcbegin{artifact}%
    \hrule height 1pt%
  \tagmcend%
  \vskip 1em%
}
\makeatother
\title{Algebra First-Year Exam, May 2014, Part 2}
\author{}
\date{}
\begin{document}
\maketitle
\thispagestyle{empty}
\begin{examinstructions}
Answer four problems. (If you turn in more, the first four will be graded.) Within reason, you may quote theorems as long as you state them clearly.
\end{examinstructions}
\begin{examproblems}{0}{H2}
\def\ProblemHeadingText{Problem 1.}
\item
Let \(R=\mathbb{Z}[i]\) be the ring of Gaussian integers. Let~\(M\) be a maximal ideal of~\(R\). Prove that \(R/M\) is a finite field of order~\(p\) or \(p^2\) where~\(p\) is a (rational) prime number.
\def\ProblemHeadingText{Problem 2.}
\item
This problem has three parts.
\begin{examsubparts}
\item[\textnormal{(a)}]
Give an example of an integral domain which is not a unique factorization domain.
\item[\textnormal{(b)}]
Give an example of a principal ideal domain which is not a field and is not isomorphic to~\(\mathbb{Z}\).
\item[\textnormal{(c)}]
Give an example of a commutative ring with identity which is not an integral domain.
\end{examsubparts}
\def\ProblemHeadingText{Problem 3.}
\item
This problem has two parts.
\begin{examsubparts}
\item[\textnormal{(a)}]
Define Euclidean domain.
\item[\textnormal{(b)}]
Prove that every Euclidean domain is a principal ideal domain.
\end{examsubparts}
\def\ProblemHeadingText{Problem 4.}
\item
Let~\(R\) be a ring with identity and let~\(M\) be an~\(R\)-module. Prove that the annihilator \(\operatorname{Ann}(M)\) is a two sided ideal of~\(R\).
\def\ProblemHeadingText{Problem 5.}
\item
Let~\(V\) be a vector space of dimension 4 over the field~\(\mathbb{Q}\) of rational numbers. Let \(T:V\to V\) be a linear transformation. Suppose that \((T^2-\operatorname{Id})^2=0\).
\begin{examsubparts}
\item[\textnormal{(a)}]
What are the possible minimal polynomials for~\(T\)?
\item[\textnormal{(b)}]
For each possible minimal polynomial from (a), give a matrix with this minimal polynomial.
\end{examsubparts}
\def\ProblemHeadingText{Problem 6.}
\item
Let \(F/K\) be an extension of fields, and suppose that \([F:K]=5\). Prove that there exists an irreducible polynomial \(p(x)\in K[x]\) such that~\(F\) is isomorphic to \(K[x]/(p(x))\).
\end{examproblems}
\end{document}
